\section*{Supersymmetric Algebras}

For $n \in \mathbb{Z}_{\geq 0}$ the standard Clifford algebra $\mathrm{Cl}_{+n}$ is the unital algebra over $\mathbb{R}$ generated by $e_1, \ldots, e_n$ subject to the relation
\[
e_i e_j + e_j e_i = 2 \delta_{ij}, \quad 1 \leq i, j \leq n.
\]

$\mathrm{Cl}_{-n}$ is similar, but the relation is $e_i e_j + e_j e_i = -2 \delta_{ij}$.
The complex version, for which the sign is irrelevant but we take the plus sign for definiteness, is denoted $\mathrm{Cl}_n^{\mathbb{C}}$.

The abstract version for a vector space $U$ (complex or real) with a symmetric bilinear form $B$ is the free unital algebra $\mathrm{Cl}(U, B)$ generated by $U$ subject to
\[
u_1 u_2 + u_2 u_1 = B(u_1, u_2), \quad u_1, u_2 \in U.
\]

These Clifford algebras are superalgebras.

\begin{question}
(a) Construct an isomorphism of superalgebras $\mathrm{Cl}_2^{\mathbb{C}} \cong \mathrm{End}(\mathbb{C}^{1|1})$.
\end{question}

\begin{question}
(b) Define the tensor product $A_1 \otimes A_2$ of superalgebras $A_1, A_2$. Mind your Koszul sign rule!
\end{question}

\begin{question}
(c) Let $(U_i, B_i)$ be real vector spaces and symmetric bilinear forms. Construct an isomorphism
\[
\mathrm{Cl}(U_1, B_1) \otimes \mathrm{Cl}(U_2, B_2) \cong \mathrm{Cl}(U_1 \oplus U_2, B_1 \oplus B_2).
\]
\end{question}

\begin{question}
(d) Construct an isomorphism of superalgebras $\mathrm{Cl}_{+8} \cong \mathrm{End}(\mathbb{R}^{8|8})$. (You may want to apply (c).)
\end{question}

\begin{question}
(e) Observe that $\mathbb{R}^n \subset \mathrm{Cl}_{+n}$. For $\xi \in \mathbb{R}^n$ a vector of unit norm, check that the map $\zeta \mapsto -\xi \zeta \xi^{-1}$ reflects $\zeta \in \mathbb{R}^n$ in the hyperplane perpendicular to $\xi$.
\end{question}

\begin{question}
(f) For $\theta \in \mathbb{R}$ define $g_\theta = \cos(\theta/2) + \sin(\theta/2) e_1 e_2$. What transformation does $g_\theta$ induce on $\mathbb{R}^n$?
\end{question}

\begin{question}
(g) Show that “quadratic” elements
\[
\frac{1}{2} A^{ij} e_i e_j
\]
in the Clifford algebra are closed under commutation. Can you identify this Lie algebra? (Here the $A^{ij}$ are scalars and the expression is summed over the repeated indices $i, j$.)
\end{question}

\section*{Symplectic Geometry and Hamiltonian Dynamics}

\subsection*{1. Let $N$ be a smooth manifold of dimension $n = 2m$, and suppose $\omega \in \Omega^2_N$.}

\begin{enumerate}
    \item[(a)] Prove that $\omega$ is nondegenerate (at each point of $N$) iff $\omega^m \in \Omega^{2m}_N$ is everywhere nonzero.

    We write for a choice of coordinates $(q^{i}, p_{i})$, such that $dq^{i} \wedge dp_{i} \neq 0$. Such a choice of coordinates exists by choosing $2n$ orthogonal vectors in the tangent space $T_{p} M$ and looking in restriction of a chart around $p$. 

    $$\omega = \sum_{i = 1}^{m} dq^{i} \wedge dp_{i}$$

    Then, 

    $$\omega^{m}/m! = dq^{1} \wedge dp_{1} \wedge dq^{2} \wedge dp_{2} \dots dq^{n} \wedge dp_{n}$$

    Therefore, $\omega^{m}/m!$ does not vanish in this chart. Now recall the fact that if a tensor vanishes in one local choice of coordinates it will vanish in all others - this follows from the tensor transformation laws. Therefore, it sufficed to show that $\omega^{m}/m!$ did not vanish in a single chart. 
    
    Now, for the other direction. We note that the right hand side is non-zero by definition. For the the other direction: note that 

    $$\omega^{m}/m! = \omega \wedge \omega \dots \omega$$

    Since we are in even dimensional space, $\wedge^{2n} M$, this vanishes only if $\omega = 0$. 

    \item[(b)] Assume now that $\omega$ is a symplectic form. Prove that the subspace $\mathfrak{X}^\omega_N \subset \mathfrak{X}_N$ of symplectic vector fields is closed under the Lie bracket.

    Follows from the identity: $$\mathcal{L}_{[X_{f_{1}}, X_{f_{2}}]} \omega = [\mathcal{L}_{X_{f_{1}}}, \mathcal{L}_{X_{f_{2}}}] \omega = 0$$
    \item[(c)] Define a Lie algebra structure on $\Omega^0_N$ by transport from the Lie bracket on $\mathfrak{X}^\omega_N$ via the symplectic gradient map $f \mapsto \xi_f$. Prove:
    \[
    \{f, g\} = \xi_f \cdot g = \omega(\xi_f, \xi_g), \quad f, g \in \Omega^0_N.
    \]

    % $$\{f, g\} = \xi_f \cdot g = \omega \left( \omega^{-1} df, \omega dg \right) = \omega(\xi_{f}, \xi_{g})$$

    \item[(d)] A \emph{Hamiltonian vector field} is a symplectic gradient of a function. Is the subspace of Hamiltonian vector fields invariant under Lie bracket?

    Use the identity:

    $$X_{\{f,g\}} = [X_{f}, X_{g}]$$

    Therefore:

    $$\{f, g \} = \omega ^{-1} [X_{f}, X_{g}]$$

    So we have shown that $[X_{f}, X_{g}]$ is again Hamiltonian,. 
    
\end{enumerate}

\subsection*{2. Let $V$ be a real symplectic vector space and $A$ an affine space over $V$.}

\begin{enumerate}
    \item[(a)] Define the space of affine polynomial functions on $A$ and prove that it is closed under the Poisson bracket.

    \item[(b)] Prove that the space of affine linear functions is also invariant. Do you recognize the resulting Lie algebra?

    \item[(c)] Prove that the space of affine quadratic functions is also invariant. Do you recognize the resulting Lie algebra?

    \item[(d)] Does the pattern continue?
\end{enumerate}

\newpage

\section*{Quantum Theory and Classical Mechanics}

\begin{enumerate}
    \item Consider a classical particle on the Euclidean line $\mathbb{E}^1$ with coordinate $x$. Fix $a > 0$ and let $V(x) = \frac{1}{2}(x^2 - a^2)^2$.
    \begin{itemize}
        \item What are the classical trajectories of least energy?


        Let $m = 1$. Then:

        $$\ddot{x} = (x^2 - a^2) \cdot 2x$$

        When the particle has least energy, then $V'(x) = 0$, so $\pm a = x$, and $V(x) = 0$. 

        $$\ddot{x} = 0, x = a$$

        So the particle is trapped at the bottom of the wells. 

        \item Repeat for the inverted potential $-V$.

        Now the minimum is when:

        $$V'(x) = 0$$

        So, $x = 0$, and we again find that the particle is trapped at the bottom of a well. 
        

        \item Given $T > 0$, is there a classical trajectory $\gamma : [-T/2, T/2] \to \mathbb{E}^1$ such that $\gamma(-T/2) = -a$ and $\gamma(T/2) = a$?

        In this case, we must have:

        $$\frac{1}{2} \dot{x}_{I}^2 = a^4/2$$

        Therefore:

        $$\frac{1}{2} \dot{x}^2 = a^4/2 - \frac{1}{2}(x^2 - a^2)^2$$

        $$dt = \frac{1}{\sqrt{a^4 - (x^2 - a^2)^2}} dx$$

        And the total time between the endpoints for the path $\gamma$ is:

        $$T = \int_{-a}^{a} \frac{1}{\sqrt{a^4 - (x^2 - a^2)^2}} dx$$

        Since this integral has a non-removable singularity at the origin, we have that $T = \infty$ and we have a contradiction. 
        
        \item What happens as $T \to \infty$?

        In this case, it becomes possible. 
    \end{itemize}

    \item[(3)] Particle on a ring $M = \mathbb{E}^1 / 2\pi \mathbb{Z}$ of length $2\pi L$.
    \begin{enumerate}
        \item[(a)] Work out the details of Example 3.57 in the notes. We leave the reader to identify the space $N$ of classical solutions and show that
        \begin{itemize}
         \item the Hamiltonian system derived from (3.57) is independent of $\theta$;
        \item the principal bundle (3.53) with connection changes with $\theta$.
        \end{itemize}

        \begin{equation}
        L = \left( \frac{m}{2} \dot{x}^2 + \frac{\theta}{2\pi} \dot{x} \right) dt.
        \end{equation}

        To find the Hamiltonian system: we do a Legrendre Transformation:

        $$p_{x} = \frac{\partial L}{\partial \dot{x}} = m \dot{x} - \frac{\theta}{2 \pi}$$

        $$p_{\theta} = \frac{\partial L}{\partial \dot{\theta}} = 0$$

        $$H = p_{x} \dot{x} - L = m \dot{x} \dot{x} -  \frac{\theta}{2 \pi} \dot{x} - L$$

        $$H = m/2 \dot{x}^2$$

        So the Hamiltonian system is independent of $\theta$. 

        Define the connection:

        $$\gamma = \theta \frac{x}{2 \pi}$$

        Then when we work with Lagrangians with fixed boundary conditions, then:

        $$\gamma(t_{1}) - \gamma(t_{0}) = \delta \int_{t_{0}}^{t_{1}} \mathcal{L} dt$$

        So, we have appropriately identified the connection in our bundle:

        $$M = N \times T$$


        \item[(b)] Quantum theory has Hilbert space $\mathcal{H} = L^2(M; \mathbb{C})$, and Hamiltonian:
        \[
        H = \frac{1}{2} \left( i \frac{\partial}{\partial x} + \frac{1}{2\pi} \theta \right)^2.
        \]
        Diagonalize $H$ and compute its spectrum. Is there periodicity in $\theta$? Compare to the classical system.


        Note that just as in the classical system:

        $$H = \frac{1}{2} p_{x}^2$$

        Where now - $p_{x} \to \hat{p_{x}}$

        $$\hat{H} = -\frac{1}{2} \frac{\partial^2}{\partial x^2} + \frac{i \theta}{2 \pi} \frac{\partial}{\partial x} + \frac{\theta^2}{8 \pi^2}$$

        Consider basis solutions of the form:

        $$\psi = e^{i \alpha x}$$

        The Schrödinger equation is:

        $$-\frac{1}{2} \frac{\partial^2 \psi}{\partial x^2} + \frac{i \theta}{2 \pi} \frac{\partial \psi}{\partial x} +  \frac{\theta^2}{8 \pi^2} = H_{\pm} \psi$$

        Note the following identity:
        $$\frac{1}{2} \alpha^2 - \frac{\theta}{2 \pi} \alpha + \frac{\theta^2}{8 \pi^2} = \frac{1}{2} (\alpha - \frac{\theta}{2 \pi})^2$$

        Therefore, we say that:

        $$E_{\alpha, \theta} = \frac{1}{2} (\alpha - \frac{\theta}{2 \pi})^2$$

        Since $\psi$ is valued on $S^{1}$, we note that $\alpha \cdot 2\pi \in 2 \pi \mathbb{N}$. Therefore:

        
        $$E_{n, \theta} = \frac{1}{2} (n - \frac{\theta}{2 \pi})^2$$

        Where $n$ is an integer. Considering the spectrum as a whole:

        $$\operatorname{Spec}(\theta) = \cup_{n} E_{n, \theta}$$

        We note that:

        $$\operatorname{Spec}(\theta) = \operatorname{Spec}(\theta \pm 2 \pi)$$

        So the spectrum is invariant under rotations as well. 
       
    \end{enumerate}

    % One approach: maximize $\langle \psi, T\psi \rangle$ over the unit ball. Use that $T$ preserves orthogonal complements to iterate the argument.
    

    \item[(4)] Let $\mathcal{H}$ be a separable complex Hilbert space. A continuous linear operator $T : \mathcal{H} \to \mathcal{H}$ is compact if the closure of $T(D) \subset \mathcal{H}$ (where $D$ is the unit disk) is compact. Prove:

    \textbf{Theorem.} Let $T$ be a positive self-adjoint compact operator. Then there exists an orthonormal basis $\{e_n\}_{n=1}^\infty$ and positive numbers $\mu_1 \ge \mu_2 \ge \cdots$ such that:
    \begin{enumerate}
        \item[(i)] $T\psi_n = \mu_n \psi_n$
        \item[(ii)] For any $c > 0$, only finitely many $\mu_n > c$
        \item[(iii)] $\lim_{n \to \infty} \mu_n = 0$
    \end{enumerate}

    We maximise $\langle \psi, T\psi \rangle$ over the unit ball. A maximum exists since $T$ is compact. Declare:

    $$\langle \psi, T \psi \rangle = \mu$$

    $$\langle \psi \vert T - \mu \vert \psi \rangle = 0$$

    Therefore:

    $$T \psi = \mu \psi$$

    Now consider the orthogonal complement of $\psi$ in $\mathcal{H}$, and run this argument again. We end up finding $\psi_{1}, \dots, \psi_{\infty}$ such that:

    $$T \psi_{n} = \mu_{n} \psi_{n}$$

    Suppose for sake of contradiction that infinitely many $\mu_{n} > c$. Then:

    $$T = \sum_{k} \mu_{n_{k}} \vert e_{n_{k}} \rangle \langle e_{n_{k}} \vert$$

    Where $e_{n_{k}}$ are the normalised basis of $\psi_{1}, \dots, \psi_{\infty}$. Now we note that: $||T|| = \infty$. Contradiction.

    As for any $c > 0$, we note there are finitely many $\mu_{n} > c$ yet there are infinitely many $\mu_{n}$, we note there is an accumulation point of $\mu_{n}$ around $0$. Therefore:

    $$\lim_{n \to \infty} \mu_{n} \to 0$$

    
\end{enumerate}

\newpage

\section*{Problem Set \#4: Principal Bundles and Descent}

\subsection*{1.}

The construction of the principal $\mathbb{R}$-bundle with connection $T \to N$ in §3.4 of the notes is clarified below.

In general, if $\pi: M \to N$ and $\mathscr{O}$ is a contravariant local object on $M$ (e.g. a differential form, bundle, etc.), we can:
\begin{itemize}
    \item Pushforward: Produce $\pi_* \mathscr{O}$ (changes geometric nature).
    \item Descent: Produce $\overline{\mathscr{O}}$ on $N$ and an isomorphism $\mathscr{O} \cong \pi^* \overline{\mathscr{O}}$.
\end{itemize}

We work in the descent context. Let:

\[
P = (N \times \mathbb{R}) \times \mathbb{R}_s \rightarrow N \times \mathbb{R}_t
\]

\begin{enumerate}
    \item[(a)] Prove that the 1-form:
    \[
    \Theta = ds + p^* \gamma \in \Omega^1_P
    \]
    is a connection. Compute the curvature $\Omega \in \Omega^2_{N \times \mathbb{R}}$.

    To show a $1$-form is a connection on a principle $G$-bundle, we must show:

    $$\theta(X^{\#}) = X, X \in \mathfrak{g}$$

    The connection eats the fundamental vector field and returns the generator.
    Souce: \href{https://empg.maths.ed.ac.uk/Activities/GT/Lect1.pdf}{here}, page $6$. 
    In our case, since the bundle is trivial, 

    $$X^{\#} = \frac{\partial}{\partial s}$$

    We note that:

    $$\frac{\partial}{\partial s} \left[ ds + p^{*} \gamma \right] = 1 \in \mathbb{R}$$

    So, $\theta$ is a connection. Now, we compute its curvature:

    $$\Omega = d \Theta = d(p^{*} \gamma) = p^{*} d \gamma$$

    Since pullbacks and exterior derivatives commute. 

    \item[(b)] Show that curvature descends under projection $\pi : N \times \mathbb{R} \to N$. Construct $\omega \in \Omega^2_N$ such that $\Omega = \pi^* \omega$.

    Consider the $\mathbb{R}$ bundle- $N \times \mathbb{R} \to \mathbb{R}$. Define $\omega = d \left[\pi^{*} \circ p^{*} \gamma \right]$. Then, $\omega \in \Omega_{N}^{2}$ such that the property holds

    \item[(c)] Descend $P \to N \times \mathbb{R}$ and its connection $\Theta$ to a principal $\mathbb{R}$-bundle $T \to N$ with connection. Hint: $T_n$ is the $\mathbb{R}$-torsor of parallel sections of $P$ over $\{n\} \times \mathbb{R}$.

    We have the bundle:

    $$(N \times \mathbb{R}) \times \mathbb{R}_{s} \to N \times \mathbb{R}_{t}$$

    We want to construct a new bundle:

    $$T \to N$$

    Out of this one - the trouble is that we have to drop $\mathbb{R}_{t}$ in the base space in our bundle. Consider $T_{n}$ the $\mathbb{R}$-torsor of parallel sections of $P$ over $\{n\} \times \mathbb{R}$. Then: we hvae a bundle:

    $$T \to N$$

    $$T = \cup_{n} T_{n}$$

    Where $\cup_{n}$ is the disjoint union. Now, the new connection is:

    $$\left[s_{1}(n, t_{1}) - s_{2}(n, t_{0}) \right] - \left[s_{1}(n, t_{0}) - s_{2}(n, t_{0}) \right]= \int_{t_{0}}^{t_{1}} \gamma(t) dt $$

    $$\left[s_{1}(n, t_{0}) - s_{2}(n, t_{0}) \right] = C$$


    \item[(d)] Consider translation action of $\mathbb{R}_t$ on $N \times \mathbb{R}_t$. Define a lift to $P$ such that:
    \begin{itemize}
        \item[(i)] $\Theta$ is invariant,
        \item[(ii)] If $\xi = \partial/\partial t$ generates the translation, then $\iota_\xi \Theta = 0$.
    \end{itemize}
    Conclude $\Theta$ descends to the quotient of $P$.

    Elements of $P$ are of the form:

    $$(n, t, s) \to (n, t)$$

    Consider the quotient of $P$, $P/\mathbb{R}_{t} \to N$, given by: $[p] \to n$. Note that:

    $$\Theta(X^{\#}) = 0$$

    Therefore, $\Theta$ is a connection on $[p] = [(n, t, s)] = (n, 0, s)$. 

    \item[(e)] More generally, suppose $\pi: Q \to N$ is a principal $G$-bundle. What descent conditions must a differential form $\Omega \in \Omega^q(Q)$ satisfy? Characterize the image of $\pi^*: \Omega^q(N) \to \Omega^q(Q)^G$.



\end{enumerate}

